\printedpage{1}
% head: A. Grothendieck : Dérivateurs, Ch. XVIII, Édition des Universités de Montpellier II et Paris VII
% foot: Version 17 décembre 2025

\msp{page 50}

\section*{4\quad Paratopos and the operation $\mathcal{M}\boxtimes\mathcal{N}$}

\subsection*{4.1\quad Pseudo-topos. Problem of the product.}

I call \textit{pseudo-$\mathfrak{U}$-topos category}, or simply \textit{pseudo-topos}, a $\mathfrak{U}$-category $\mathcal{M}$ stable under $\mathfrak{U}$-direct limits, and which admits a \textit{small} subcategory $\mathcal{M}_0$ that is \textit{strictly generating} (i.e. generating by strict epimorphisms). Then $\mathcal{M}$ also admits $\mathfrak{U}$-inverse limits (SGA 4 I 8.12.7). (But the example of (Gr), cf. loc. cit. 8.12.9, shows that $\mathcal{M}^\circ$ is not necessarily, for all that, a pseudo-$\mathfrak{U}$-topos, for lack of $\mathcal{M}$ admitting a small strictly \textit{co}generating subcategory. But theorem 2.6 \texttt{[of chapter XIX]} gives conditions on $\mathcal{M}$ (mostly of the exactness type) % [sic] par forcément -> pas forcément, misprint; also the parenthesis opened at "(Mais" is never closed
\msp{page 51}
for $\mathcal{M}$ to admit enough injectives -- which will then imply the existence of a small strictly cogenerating subcategory, hence that $\mathcal{M}^\circ$ is also a pseudo-topos (\footnote{This is the case in particular if $\mathcal{M}$ is a topos.}).

I visualize the pseudo-topos categories as being the category of ``sheaves'' on an object $\mathcal{X}$, visualized as a kind of ``space'' which I distinguish from the category $\mathcal{M}=\Fais(\mathcal{X})$ of sheaves on $\mathcal{X}$. It is $\mathcal{X}$ which will, for this, be the ``pseudo-topos''. The pseudo-topoi form a 2-category, the variance in $\mathcal{X}$ being opposite to the one in $\Fais(\mathcal{X})=\mathcal{M}$. To be precise, if $\mathcal{X}$ corresponds to $\mathcal{M}$, $\mathcal{Y}$ to $\mathcal{N}$, a morphism $f$ from $\mathcal{X}$ to $\mathcal{Y}$ will be given by a functor
\[
f^* \;:\; \mathcal{N} \longrightarrow \mathcal{M}
\tag*{(4.1.1)}
\]
in the opposite direction, commuting with arbitrary direct $\varinjlim$ (\footnote{Say here that the objects of $\mathcal{M}$ are visualized as generalized ``open sets'' of $\mathcal{X}$. We have $\mathcal{M}\xrightarrow{\sim}\uHom^!(\mathcal{M}^\circ,\Ens)$, i.e. $\mathcal{M}$ is well and truly interpreted as the ``sheaves'' on this category $\mathcal{M}$ of open sets \ldots}). (\textbf{N.B.} if $\mathcal{X}$, $\mathcal{Y}$ are even topoi, i.e. $\mathcal{M}$, $\mathcal{N}$ topos-categories, and if we want $f$ to be a morphism of topoi, we will require moreover that $f^*$ be moreover \textit{left exact}. We will say that a \textit{morphism of pseudo-topoi is exact} if moreover $f^*$ is left exact.)

\msp{page 52}
By virtue of the dual of loc. cit. 8.12.7, this is equivalent to the fact that $f^*$ admits a right adjoint $f_*$,
\[
f_* \;:\; \mathcal{M} \longrightarrow \mathcal{N}\;,
\tag*{(4.1.2)}
\]
which therefore commutes with small $\varprojlim$\,. The datum of $f_*$ satisfying this property therefore determines $f^*$ up to unique isomorphism. But in general, it is not enough to give oneself such an $f_*$, in order to be able to assert that it corresponds to an $f^*$, i.e. that it admits a right adjoint\,; for this, we would need to know that the functors $\mathcal{M}\longrightarrow\Ens$ which commute with small $\varprojlim$ are representable, which will be the case if $\mathcal{M}$ admits a small \textit{co}generating subcategory, i.e. if $\mathcal{M}^\circ$ is also a pseudo-topos category. % [sic] qu'il admette un adjoint à droite -> adjoint à gauche, f^* is the left adjoint of f_*

\newpage